\section{问题四：波动电价下的购电调度}
\subsection{4-2：价格与净负载联合场景}
题面给出了实际价格，但未明确每天0点是否已知全天价格。因此主模型采用“当天价格未知、仅有历史观测”的信息条件；另设当日电价已知的对照。实际价格用于事后结算，不能因附件包含全年价格而提前使用未来值。

以1月8—31日的24个历史预测日比较前一日、近7日均值、近7日形状加前一日水平、最近4次同星期几均值四种规则，按预测RMSE选定最后一种，并从2月起固定。样本日期必须早于决策日，预测次日时也不能使用尚未实现的当日价格。缺少同星期样本时用近7日均值回退。

将与净负载残差同日的电价预测误差配对，得到
\begin{equation}
p_{s,t}=\max\{0.001,\widehat p_{d,t}+p_{r_s,t}-\widehat p_{r_s,t}\},\qquad
\overline p_{d,t}=\sum_s\omega_s p_{s,t}.
\end{equation}
0.001元/kWh为场景正值下限，不改变原始实际价格。模型保持第二问的共享日前控制、56天等权及48小时前瞻，目标改为
\begin{equation}
\min\ \sum_t\overline p_{d,t}g_t+
5\sum_s\omega_s\sum_t p_{s,t}e_{s,t}+
\sum_t\widehat p_{d+1|d,t}g_{144+t}.\label{eq:q42}
\end{equation}
对于净供给$z_t=g_t+b_t-c_t$，一般有
$\E[p_t(N_t-z_t)_+]\ne\E[p_t]\E[(N_t-z_t)_+]$，故价格与缺口的依赖关系可能影响风险权重。这里直接相关的是同一时段的联合分布，保留整日配对本身不能保证额外收益。均价对照使用相同净负载场景，将所有价格场景替换为其均值，以检验该关系的实际作用。

\begin{table}[htbp]\centering\small
\caption{4-2同价结算下的价格信息对照}
\begin{tabular}{lrr}\toprule
方案 & 实际费用/元 & 较重结算基线节省/元\\\midrule
原Q2计划重结算 & 15184704.23 & 0.00\\
均价对照 & 15160921.45 & 23782.79\\
联合场景主方案 & 15163043.89 & 21660.34\\
当日电价已知对照 & 15033043.18 & 151661.05\\
\bottomrule\end{tabular}\end{table}

\fig{07_q42_price_information}
联合场景主方案费用15163043.89元，相对第二问原计划按附件4实际价格重结算的15184704.23元节省21660.34元，约0.143\%。均价对照费用15160921.45元，反而比联合方案低2122.44元。因此，本年结果支持存在计划费与紧急费的权衡，不能据此声称联合场景必然更优。当日电价已知对照为15033043.18元，放宽了信息条件，但仍受既有预测和滚动近似限制，不是已证明的理论最优下界。

\fig{08_q42_schedule_shift}
图\ref{fig:08_q42_schedule_shift}对齐实际电价、场景均价与主方案相对重结算基线的日前购电变化。两条价格曲线间的填色是差值提示，不是置信区间。购电变化还受到净负载场景、储能状态和次日前瞻影响，不能把每个变化都归因于相邻时段的价格偏差。
\FloatBarrier

\subsection{4-3：随机价格下的日内调整与现有实现边界}
沿用第三问的融合、场景树及调整合同，未来电价预测采用此前7天同一时段均值，已执行时段使用已知价格。电价误差与净负载误差来自相同历史日，场景正值下限为$10^{-5}$元/kWh。其价格预测规则与4-2不同，应分别说明，不能合写成同一个电价预测器。

\subsubsection{优化目标与实际结算的区别}\label{sec:q43objective}
现有主程序除0点原计划价格项、增购/取消项、紧急补救项及次日前瞻外，还给节点调整后购电$q_{n,t}$设置了价格系数。令$\widetilde p_t$为当前历史价格预测，$\overline p_{n,t}$为节点内价格场景均值，则0点实际求解的目标可表示为
\begin{align}
\widetilde J_{43}={}&\sum_t\widetilde p_t g_t^0
+\sum_n\pi_n\sum_{t\in\mathcal T_n}\overline p_{n,t}q_{n,t}\notag\\
&+\sum_n\pi_n\sum_{t\in\mathcal T_n}\overline p_{n,t}(1.5u_{n,t}-0.5v_{n,t})\notag\\
&+\frac5S\sum_s\sum_t p_{s,t}e_{s,t}+\widetilde\Phi.\label{eq:q43actual}
\end{align}
后续更新时原计划固定，其首项省略为常数。现有次日前瞻使用当前价格预测向量，并按对应节点概率加权；它是原实现的近似，非次日价格的完美信息。

式\eqref{eq:q43actual}中的$\sum\pi_n\overline p_{n,t}q_{n,t}$在第三问优化中不存在，也不属于既定合同实际费用的额外应付款。实际报告仍按式\eqref{eq:q3settle}代入附件4价格计算。因而当前4-3结果是这一既有决策目标生成、按合同重新结算的策略结果，不能表述为对合同期望总费用目标的完全一致求解。该附加项可能影响购电和库存，尚不能在未重新优化对照的情况下量化其影响；本文保留正式结果并将此列为实质性待核对问题，不为附加项补造未经记录的经济动机。

\subsubsection{实际费用及与第三问的差异}
\begin{table}[htbp]\centering\small
\caption{4-3主结果实际费用分解}
\begin{tabular}{lr}\toprule
项目 & 费用/元\\\midrule
原计划费 & 13457230.68\\
增购费 & 152250.38\\
取消退款（扣减） & -173200.59\\
违约费 & 86600.30\\
净调整费 & 65650.08\\
紧急费 & 1538854.67\\
总费用 & 15061735.43\\
\bottomrule\end{tabular}\end{table}

4-3实际总费用15061735.43元，由原计划费13457230.68元、净调整费65650.08元和紧急费1538854.67元组成。相对第三问增加1039338.45元，其中紧急费增加877139.21元，占账面差额84.39\%。该比例只说明费用分项，不证明额外成本全部由电价波动或预报误差造成。

\fig{09_q43_incremental_comparison}
图\ref{fig:09_q43_incremental_comparison}同时区分跨情形费用差额与情形内部新增更新的效果。4-3从本问历史预测对照到四次更新，共减少455249.11元；其中18点相对0、6、12点方案减少6875.04元。它与第三问18点增加38.88元的结果不同，进一步说明更新时点价值依赖价格、预测与决策目标，不能把固定电价下的频率结论直接移植过来。

\begin{table}[htbp]\centering\small
\caption{4-3更新频率对照}
\begin{tabular}{lrr}\toprule
方案 & 总费用/元 & 紧急购电量/kWh\\\midrule
同口径历史预测 & 15516984.54 & 504526.13\\
仅0点融合 & 15295120.10 & 461933.56\\
0、6点 & 15173094.46 & 410511.06\\
0、6、12点 & 15068610.47 & 387077.17\\
四次更新主方案 & 15061735.43 & 383981.09\\
\bottomrule\end{tabular}\end{table}

\fig{10_q43_schedule_differences}
图\ref{fig:10_q43_schedule_differences}按同日同段将4-3减第三问，分别展示价格、最终购电功率和段末储电量差值。每个面板保留四日共576个十分钟值，使用以零为中心的全范围对称色标。两问年度初末状态相同，但指定日初库存可能已不同，因此库存差含前期历史路径影响。此处还有式\eqref{eq:q43actual}所示目标差异，图中共同变化不能用作单一价格因素的因果识别。
\FloatBarrier
